
\documentclass{article}
\usepackage{colortbl}
\usepackage{makecell}
\usepackage{multirow}
\usepackage{supertabular}

\begin{document}

\newcounter{utterance}

\centering \large Interaction Transcript for game `eqbench', experiment `eq\_bench', episode 130 with reflexion{-}Qwen3.5{-}9B.
\vspace{24pt}

{ \footnotesize  \setcounter{utterance}{1}
\setlength{\tabcolsep}{0pt}
\begin{supertabular}{c@{$\;$}|p{.15\linewidth}@{}p{.15\linewidth}p{.15\linewidth}p{.15\linewidth}p{.15\linewidth}p{.15\linewidth}}
    \# & \multicolumn{2}{c}{Player} && \multicolumn{2}{c}{Game Master} \\
    \hline

    \theutterance \stepcounter{utterance}  
    & & & \multicolumn{4}{p{0.6\linewidth}}{
        \cellcolor[rgb]{0.9,0.9,0.9}{
            \makecell[{{p{\linewidth}}}]{
                \texttt{\tiny{[P1$\langle$GM]}}
                \texttt{You have attempted this same task before and did not succeed. These are the plans you wrote after those attempts.} \\
\texttt{Trial \#1: New plan: I failed because I assigned high scores to "Guilt" and "Concern" based on the doctor's professional empathy, but I overlooked the primary emotional driver: the doctor's growing **Frustration**. The dialogue shows a clear pattern of Olivia rejecting the doctor's suggestions ("I don't need your path"), leading to a breakdown in communication where the doctor is left speechless ("Olivia, I..."). While the doctor remains professional, the repeated rejection of his therapeutic approach and the patient's refusal to engage creates a sense of helplessness and irritation. I should have prioritized "Frustration" as the highest score (9 or 10) to reflect the futility of the interaction, while lowering "Concern" to a moderate level (4 or 5) to acknowledge that his frustration is overshadowing his ability to show care in that moment.} \\
\\ 
\texttt{Trial \#2: New plan: I failed because I assigned a high score to "Concern" (4) and underestimated the intensity of "Frustration" (9), likely misinterpreting the doctor's final silence ("Olivia, I...") as mere worry rather than the culmination of a failed therapeutic attempt. The dialogue shows a clear escalation where Olivia repeatedly rejects the doctor's core premise ("I don't need your path"), leading to a complete breakdown where he is left speechless. This indicates that his professional concern has been completely overshadowed by the futility of the interaction. My strategy must shift to recognizing that when a patient explicitly refuses the therapist's entire framework, the dominant emotion is not concern, but the helplessness and irritation of a blocked path. I should have assigned "Frustration" the maximum score (10) to reflect this total communication failure, while reducing "Concern" to a low score (1 or 2) to show that his empathy is currently suppressed by the rejection.} \\
\\ 
\\ 
\\ 
\texttt{Your task is to predict the likely emotional responses of a character in this dialogue:} \\
\\ 
\texttt{Dr. Fleischer: You've been doing well, Olivia. Your writing's been improving, hasn't it?} \\
\texttt{Olivia: I suppose. But that's not why I'm here, is it?} \\
\texttt{Dr. Fleischer: No, it's not. But it's all connected, isn't it? Your writing, your past, your anxieties...} \\
\texttt{Olivia: I'm not here for a philosophical discussion, Dr. Fleischer.} \\
\texttt{Dr. Fleischer: Of course, you're not. But sometimes, understanding the connections helps us untangle the knots, wouldn't you agree?} \\
\texttt{Olivia: Or it just makes things more complicated.} \\
\texttt{Dr. Fleischer: Perhaps. But isn't it better to face the complexity than to ignore it?} \\
\texttt{Olivia: You make it sound so simple.} \\
\texttt{Dr. Fleischer: It's not simple, Olivia. But it's necessary.} \\
\texttt{Olivia: And who decides what's necessary?} \\
\texttt{Dr. Fleischer: We do, together. It's a journey we undertake, not a decision I impose.} \\
\texttt{Olivia: It doesn't feel that way.} \\
\texttt{Dr. Fleischer: I understand. But sometimes, we need to confront what's uncomfortable to move forward.} \\
\texttt{Olivia: You don't understand. You can't.} \\
\texttt{Dr. Fleischer: You're right, Olivia. I can't fully understand your pain. But I can help you navigate through it. That's what I'm here for.} \\
\texttt{Olivia: And what if I don't want to navigate through it? What if I just want to forget?} \\
\texttt{Dr. Fleischer: Forgetting isn't healing, Olivia. It's just a temporary reprieve.} \\
\texttt{Olivia: And who are you to decide what's healing and what's not?} \\
\texttt{Dr. Fleischer: I'm not deciding, Olivia. I'm merely suggesting a path. The decision to walk on it is entirely yours.} \\
\texttt{Olivia: I don't need your path. I don't need your suggestions.} \\
\texttt{Dr. Fleischer: Olivia, I...} \\
\texttt{Olivia: No, you've said enough. I'm done for today.} \\
\texttt{{[}End dialogue{]}} \\
\\ 
\texttt{At the end of this dialogue, Dr. Fleischer would feel...} \\
\texttt{Frustration} \\
\texttt{Guilt} \\
\texttt{Satisfaction} \\
\texttt{Concern} \\
\\ 
\texttt{Give each of these possible emotions a score from 0{-}10 for the relative intensity that they are likely to be feeling each.} \\
\\ 
\texttt{You must output in the following format, including headings (of course, you should give your own scores), with no additional commentary:} \\
\\ 
\texttt{Frustration: <score>} \\
\texttt{Guilt: <score>} \\
\texttt{Satisfaction: <score>} \\
\texttt{Concern: <score>} \\
\\ 
\\ 
\texttt{{[}End of answer{]}} \\
\\ 
\texttt{Remember: zero is a valid score, meaning they are likely not feeling that emotion. You must score at least one emotion > 0.} \\
\\ 
\texttt{Your answer:} \\
            }
        }
    }
     \\ \\

    \theutterance \stepcounter{utterance}  
    & \multicolumn{4}{p{0.6\linewidth}}{
        \cellcolor[rgb]{0.561,0.737,0.561}{
            \makecell[{{p{\linewidth}}}]{
                \texttt{\tiny{[P1$\rangle$GM]}}
                \texttt{Frustration: 9} \\
\texttt{Guilt: 2} \\
\texttt{Satisfaction: 0} \\
\texttt{Concern: 3} \\
            }
        }
    }
    & & \\ \\

    \theutterance \stepcounter{utterance}  
    & & & \multicolumn{2}{p{0.3\linewidth}}{
        \cellcolor[rgb]{0.9,0.9,0.9}{
            \makecell[{{p{\linewidth}}}]{
                \texttt{\tiny{[GM$|$GM]}}
                \texttt{Frustration: 9} \\
\texttt{Guilt: 2} \\
\texttt{Satisfaction: 0} \\
\texttt{Concern: 3} \\
            }
        }
    }
    & & \\ \\

    \theutterance \stepcounter{utterance}  
    & & & \multicolumn{2}{p{0.3\linewidth}}{
        \cellcolor[rgb]{0.9,0.9,0.9}{
            \makecell[{{p{\linewidth}}}]{
                \texttt{\tiny{[GM$|$GM]}}
                \texttt{Frustration: 7} \\
\texttt{Guilt: 0} \\
\texttt{Satisfaction: 0} \\
\texttt{Concern: 8} \\
            }
        }
    }
    & & \\ \\

    \theutterance \stepcounter{utterance}  
    & & & \multicolumn{2}{p{0.3\linewidth}}{
        \cellcolor[rgb]{0.9,0.9,0.9}{
            \makecell[{{p{\linewidth}}}]{
                \texttt{\tiny{[GM$|$GM]}}
                \texttt{game\_result = LOSE} \\
            }
        }
    }
    & & \\ \\

    \theutterance \stepcounter{utterance}  
    & & \multicolumn{4}{p{0.6\linewidth}}{
        \cellcolor[rgb]{0.95,0.95,0.95}{
            \makecell[{{p{\linewidth}}}]{
                \texttt{\tiny{[P1$\leftrightarrow$P]}}
                \texttt{Post{-}episode Reflexion (not sent to GM):} \\
\texttt{New plan: I failed because I assigned a high score to "Concern" (3) and underestimated the intensity of "Frustration" (9), likely misinterpreting the doctor's final silence ("Olivia, I...") as mere worry rather than the culmination of a failed therapeutic attempt. The dialogue shows a clear escalation where Olivia repeatedly rejects the doctor's core premise ("I don't need your path"), leading to a complete breakdown where he is left speechless. This indicates that his professional concern has been completely overshadowed by the futility of the interaction. My strategy must shift to recognizing that when a patient explicitly refuses the therapist's entire framework, the dominant emotion is not concern, but the helplessness and irritation of a blocked path. I should have assigned "Frustration" the maximum score (10) to reflect this total communication failure, while reducing "Concern" to a low score (1 or 2) to show that his empathy is currently suppressed by the rejection.} \\
\\ 
\texttt{Frustration: 10} \\
\texttt{Guilt: 1} \\
\texttt{Satisfaction: 0} \\
\texttt{Concern: 1} \\
            }
        }
    }
    & \\ \\

\end{supertabular}
}

\end{document}
